\subsection{Interpretability and Evaluation as Instruments}
\label{sec:4.2.7}
Interpretability used to be commentary on a finished model. It is now applied during construction, at production scale. Adly Templeton and coauthors reported in 2024 that sparse autoencoders trained with up to thirty-four million features on a deployed model's middle-layer activations recover features corresponding to recognizable concepts, and that those features can be used to steer the model's behavior \autocite{templeton2024scaling}. Evaluation has moved the same way: Ethan Perez and coauthors showed in 2022 that one language model can be used to generate test cases against another, turning up tens of thousands of offensive replies from a target system without a human writing the prompts \autocite{perez2022red}.

Both are improvements on the alternative, which is shipping and finding out. What they do to this book's argument is less comfortable. The positive result reported earlier was that a commitment realized as a standing bias on what a predictive loop treats as an error worth correcting is not a component that can be located and removed. Section~\ref{sec:4.3} builds on the same property: a disposition acquired by exploration is written down nowhere, so there is no corresponding parameter to reach for. Interpretability research is the project of making such things locatable, and the steering result is the demonstration that locating them is the same operation as manipulating them. A feature that can be found and turned up can be found and turned down, by whoever holds the weights.

So the instruments cut the way the oversight schemes do, and for the same reason. An outside auditor cannot check a commitment that nobody can locate, which means the research that threatens the floor's obscurity is also the only research that could ever verify the floor exists. That is section~\ref{sec:3.7}'s trade arriving in a second place: legible memory can be tampered with, and illegible memory cannot be audited. Section~\ref{sec:11.1} carries what is known about which side is winning: published defenses against tampering are followed by attacks built against them, which yields a price curve rather than a guarantee.
